\section{The earlier editions}\label{sec:earlier}

\subsection{The continuations of 14--21 September}

Their gap and heat bounds are superseded by Sections~\ref{sec:gap} and~\ref{sec:heat}; their other calculations remain in the record as dated editions.
\begin{itemize}
\item \emph{14--16 September} (the 297-page web-continuation reader). The first gap theorems of the same architecture:
\begin{itemize}
\item a gauge-native order-1 theorem (equations G22 and G26--G29), with $\Delta_L>2.07445\kappa$ at $g^2\ge5$;
\item a single-norm order-2 theorem (R10), for $g^2\ge\sqrt{(32+\sqrt{354})/3}\approx4.1156$;
\item an order-2 theorem (L17) for $g^2\ge3.8260$, which uses a computed bound $944984/351$ for the cubic term of the residual.
\end{itemize}
\status{Located; the proofs were not read.
\begin{itemize}
\item The order-1 theorem is Corollary~\ref{cor:ym-orderN}(a), proved there.
\item R10 is strengthened by Corollary~\ref{cor:ym-orderN}(b): the threshold is $4.0187$ with the workbench's order-2 pair, and $3.9781$ with the first pass's order-2 values.
\item L17's computed bound was not checked.
\item One arithmetic line of this edition is wrong. It justifies the self bound $40/27$ of its table (C21) by ``$38/81+1564/810=40/27$'' (p.~153), but the left side equals $12/5$. The value $40/27=3\cdot\frac1{81}\cdot8+15\cdot\frac1{810}\cdot64$ is the one that the referee pass of this reader obtained for the self cluster, so the bound itself stands (found by the referee pass; the arithmetic checked here).
\end{itemize}}
\item \emph{17 September} (the 49-page quartic cube reader). The fourth-order vacuum source in two presentations, equal as polynomials (78 representatives, 8{,}621 anchored multisets); the sixth-order vacuum energy, with the contribution $-83/1944$ of the elementary cube; and the gap bound $\Delta_L>1.6584\kappa$ for $g^2\ge15/4$.
\status{The gap bound is Corollary~\ref{cor:ym-orderN}(c) at $\xi=4/225$, conditional on the order-3 and order-4 inputs; Theorem~\ref{thm:ym01} gives $d_5(4/225)>1.6993$ there. The rest located. The workbench records that this edition keeps its earlier review qualification and that the 21 September heat proof does not recertify the earlier interval (\texttt{yang-mills/README.md}).}
\item \emph{18--19 September} (inside the cumulative archives). The sixth vacuum source, the eighth-order energy, and a finite response estimate on the $240$ faces of the box $L=2$.
\status{Located. The workbench says that the new spatial tables and the final execution package were not supplied, so their completeness is reported by the source, not replayed.}
\item \emph{21 September, heat and volume} (the 46-page reader). The fourth-order heat coefficients (84 time functions over 17 geometric cases); a box-independent heat remainder on the circle $|\xi|<3/256$; the fixed-spacing spatial limit; and the full-complement energy retained above $999/1000$ at $g^2\ge16$.
\status{Superseded quantitatively by Section~\ref{sec:heat}: radius $1/55>3/256$, and a loss below $1/25000$ against $1/1000$. This edition also has a second remainder bound, U29b (p.~37), on the circle $45/4096$ with prefactor $6184/25$ and rate $15/8$. It is superseded by the family of radii of Section~\ref{sec:heat}, not by Theorem~\ref{thm:heat} alone. The five proof manuscripts are located. One of them receives the zeta programme's Gram-sandwich lemma; its steps T1--T3 were checked in \note{21\_} (Section~\ref{sec:overlaps}).}
\end{itemize}

\subsection{The foundations of 8--9 September}

The three manuscript lines of §\ref{ssec:lines} contain the workbench's first results. They are stated for the full finite Hamiltonian, with its scalar constant, open boundaries and all gauge constraints, and they say what they do not prove (§\ref{ssec:scope}).

\begin{itemize}
\item \emph{Quantum line} (79 pages).
\begin{itemize}
\item Exact extensive blocking maps, and volume-independent lower bounds on the high-energy spectral weight of selected dressed Wilson-loop channels, at every positive coupling (§11).
\item A first volume-uniform gap, $\Delta\ge\kappa(3-512\xi)$ for $0<\xi\le1/49152$, uniform in the box and the block size (§10, equation (77)). It is proved there with the creation method of Yarotsky \cite{Yarotsky}, with Baez's spin-network framework \cite{Baez} for the physical space (p.~19). The spatial line uses it for its fixed-coupling results (spatial §28: $g_0^4\ge12288$, which is the same condition).
\item A nonzero non-Abelian vertex on the box $L=2$ (§17).
\item A seven-dimensional lowest-band map that is a bijection for small coupling $g$, at a fixed box.
\item The Jacobian-map and Navier--Stokes branches (Section~\ref{sec:side}).
\end{itemize}
\item \emph{Spatial line} (129 pages). Sequences in which the box, the spacing and the coupling vary together; its conclusion (§28) states what they prove.
\begin{itemize}
\item Along a fixed-coupling diagonal with $g_0^4\ge12288$, the magnetic excitation probability eventually has zero weight on every bounded energy interval.
\item For fixed Wilson-loop states, a positive high-energy mass and a failure of strong continuity.
\item On a selected weak-coupling diagonal, the finite gaps close while the corrected magnetic probability escapes to high energy.
\item The collapse of the finite-box analyticity disk (§27; Proposition~\ref{prop:neumann}).
\end{itemize}
\item \emph{Volume line} (81 pages). Volume-independent local estimates for the actual vacuum. They include:
\begin{itemize}
\item local electric moments at every coupling;
\item the exact free physical threshold $3\kappa$ (§15);
\item estimates on the $S^6$-derived cusp at $0<\xi\le10^{-16}$;
\item §26: along the logarithmic coupling path, the upper end of its certified spectral window diverges. The document calls this a limitation of its own estimate, not a statement about the spectrum.
\end{itemize}
\end{itemize}
\status{Located, apart from the items proved or audited elsewhere in this reader: §10's gap (Proposition~\ref{prop:firstgap}), §27 (Proposition~\ref{prop:neumann}), and Theorems 13.1 and 14.2 of the quantum line (Section~\ref{sec:side}).}

\begin{proposition}[the first volume-uniform gap is implied]\label{prop:firstgap}
For $0<\xi\le3/256$ the order-1 bound of Corollary~\ref{cor:ym-orderN}(a) satisfies $d_1(\xi)\ge3-128\xi$. So the quantum line's $\Delta\ge\kappa(3-512\xi)$ on $0<\xi\le1/49152$ follows from Corollary~\ref{cor:ym-orderN}(a), which holds on the larger range $\xi\le3/256$ ($g^2\ge8/\sqrt3$).
\end{proposition}
\status{Proved here.}
\begin{proof}
With $y=256\xi/3\in(0,1]$, $d_1(\xi)=\frac32(1+\sqrt{1-y})\ge\frac32(2-y)=3-128\xi$, since $\sqrt s\ge s$ for $0\le s\le1$. And $1/49152<3/256$.
\end{proof}
